%%	Dokument-Wrapper nach der Vorlage von Herrn Plate.

%% Standard-Vorspann
\documentclass[a4paper,12pt,oneside]{report}
\usepackage[english,ngerman]{babel}                   
\usepackage[utf8]{inputenc}
\usepackage[T1]{fontenc}
\usepackage[fleqn]{amsmath}


%% Auf Schriftart Palatino umschalten
\usepackage{mathpazo}
\usepackage[scaled=.95]{helvet}
\usepackage{courier}

%% fast immer benoetigte Pakete
\usepackage{epsfig}
\usepackage{amssymb}
\usepackage{alltt}
\usepackage{latexsym}
\usepackage{makeidx}
\usepackage{textcomp}
\usepackage{rotating}
\usepackage{color}
\usepackage{caption}
\usepackage{verbatim}
\usepackage{fancyhdr}

%% spezielles Zeug
\usepackage[tablename=Tbl.]{caption}
\usepackage{nomencl}
\usepackage{tocloft}
\usepackage{Abschlussarbeit}
\usepackage{geometry}
\usepackage{graphicx}
\usepackage[onehalfspacing]{setspace}

\graphicspath{{graphics/}}
\geometry{left = 4cm, right = 2cm, top = 2cm, bottom = 2cm}

\captionsetup[table]{labelfont=bf,textfont=it}
\captionsetup[figure]{labelfont=bf,textfont=it}
\captionsetup{format=hang,indention=0pt,justification=raggedright, singlelinecheck=false}
\addto\captionsngerman{\renewcommand{\figurename}{Abb.}}

\renewcommand{\nomname}{Abkürzungsverzeichnis}

\setlength{\cftfignumwidth}{3em}

%% Neue umgebung definieren für 2 Abstracts auf einer Seite
\newenvironment{abstractDuo}{
	\begin{center}%
		\bfseries\abstractname
	\end{center}}%

%% Index erzeugen
\makeindex

%% jetzt geht's los
\begin{document}
	\raggedbottom
	
	%% Die Datei "Vorspann.tex"  bitte ausfüllen!
	\pagenumbering{Roman}
\begin{titlepage}
	\begin{figure}[h]
		 
		\begin{minipage}{0.45\linewidth} 
			\raggedright 
			\includegraphics[scale=0.15]{HMLogo.png} 
		\end{minipage} 
		\begin{minipage}{0.45\linewidth} 
			\raggedleft
			\vspace*{1 cm} 
			\includegraphics[scale=0.4]{sielogopetrolcmyk.jpg} 
		\end{minipage} 
	\end{figure}
	\begin{center} 
		\vspace{1.5 cm} {\LARGE Hochschule München} \\ 
		\vspace{1 cm} {\LARGE Fakultät für Elektrotechnik und Informationstechnik} \\ 
		\vspace{2 cm} {\Large Studiengang Elektrotechnik und Informationstechnik} \\ 
		\vspace*{2 cm}{\Huge \bf Industrie 4.0 - Pathfinding auf einer SPS\\} 
		\vspace{2 cm}
		{\bf Bachelorarbeit von Niels Garibaldi}\\
	\end{center}
	\vspace{4 cm}
	\begin{center}
		
		\begin{tabular}{ll}
			
			{\bf Bearbeitungsbeginn:} & xx.xx.2016\\
			{\bf Abgabetermin:} & xx.xx.2016\\
			{\bf lfd. Nr.:} & xxxxx\\
		\end{tabular}
	\end{center}

\end{titlepage} 

\begin{titlepage}
	\begin{figure}[h]
		
		\begin{minipage}{0.45\linewidth} 
			\raggedright 
			\includegraphics[scale=0.15]{HMLogo.png} 
		\end{minipage} 
		\begin{minipage}{0.45\linewidth} 
			\raggedleft
			\vspace*{1 cm} 
			\includegraphics[scale=0.4]{sielogopetrolcmyk.jpg} 
		\end{minipage} 
	\end{figure}
	\begin{center} 
		\vspace{1.5 cm} {\LARGE Hochschule München} \\ 
		\vspace{1 cm} {\LARGE Fakultät für Elektrotechnik und Informationstechnik} \\ 
		\vspace{2 cm} {\Large Studiengang Elektrotechnik und Informationstechnik} \\ 
		\vspace*{2 cm}{\Huge \bf Industrie 4.0 - Pathfinding auf einer SPS\\} 
		\vspace{1 cm}
		\hrule
		\vspace*{1 cm}{\Huge \bf Industry 4.0 - Pathfinding on a PLC\\} 
		\vspace{2 cm}
		{\bf Bachelorarbeit von Niels Garibaldi}\\
		\vspace{1 cm}
		{\bf betreut von  Prof. Dr. K. Ressel}\\
		
	\end{center}
	\vspace{1 cm}
	\begin{center}
		
		\begin{tabular}{ll}
			
			{\bf Bearbeitungsbeginn:} & xx.xx.2016\\
			{\bf Abgabetermin:} & xx.xx.2016\\
			{\bf lfd. Nr.:} & xxxxx\\
		\end{tabular}
	\end{center}
	
\end{titlepage}
\pagestyle{headings}
\clearpage
\phantomsection
\addcontentsline{toc}{section}{Erklärung des Bearbeiters}
\section*{Erklärung des Bearbeiters}



\begin{enumerate}
	\vspace{2 cm}
	\item
		Ich erkläre hiermit, dass ich die vorliegende Bachelorarbeit selbständig
		verfasst und noch nicht anderweitig zu Prüfungszwecken vorgelegt habe.\\
		
		Sämtliche benutzte Quellen und Hilfsmittel sind angegeben, wörtliche
		und sinngemäße Zitate sind als solche gekennzeichnet.\\
	
	\item
		Ich erkläre mein Einverständnis, dass die von mir erstellte Bachelorarbeit in die Bibliothek
		der Hochschule München eingestellt wird. Ich wurde darauf hingewiesen, dass die
		Hochschule in keiner Weise für die missbräuchliche Verwendung von Inhalten durch Dritte
		infolge der Lektüre der Arbeit haftet. Insbesondere ist mir bewusst, dass ich für die
		Anmeldung von Patenten, Warenzeichen oder Geschmacksmustern selbst verantwortlich
		bin und daraus resultierende Ansprüche selbst verfolgen muss. 
	
	\vspace{6 cm}
\end{enumerate}
\begin{flushright}
	\begin{tabular}{cl}
		München, den xx.xx.2016, & \_\_\_\_\_\_\_\_\_\_\_\_\_\_\_\_\_\_\_\_\_\_\_\_\_\_\_\_\_\_\_\_\_\_\_\_\_\_\tabularnewline
		& Niels Garibaldi\tabularnewline
	\end{tabular}
	\par
\end{flushright}
\clearpage

%% Kurzfassung/Abstract
\phantomsection
\addcontentsline{toc}{section}{Zusammenfassung/Abstract}	
\begin{abstractDuo}
	Zusammenfassungstext hier einfügen
\end{abstractDuo}
\selectlanguage{english}
\vspace{8 cm}
\begin{abstractDuo}
	Insert abstract text here
\end{abstractDuo}
\selectlanguage{ngerman}
\clearpage
%% Inhaltsverzeichnis
\pagestyle{plain}
\clearpage
\phantomsection
\addcontentsline{toc}{section}{\contentsname}

	\begin{spacing}{1.1}
	
		\tableofcontents
	\end{spacing}
\clearpage
\addcontentsline{toc}{section}{Abkürzungsverzeichnis}	
\section*{Abkürzungsverzeichnis}
\begin{acronym}[Bash]
	
	\acro{AWL}{Anweisungsliste}
	\acro{CPS}{Cyber-physisches System}
	\acro{CPPS}{Cyber-physisches Produktionssystem}
	\acrodefplural{CPPS}{Cyber-physische Produktionssysteme}
	\acro{CT}{Corporate Technologies}
	\acro{D*}{Focussed Dynamic A*}
	\acro{DB}{Datenbaustein}
	\acro{FB}{Funktionsbaustein}
	\acro{FC}{Funktion}
	\acro{FTF}{Fahrerloses Transportfahrzeug}
	\acro{FTS}{Fahrerloses Transportsystem}
	\acro{FUP}{Funktionsplan}
	\acro{IoT}{Internet of Things}
	\acro{KOP}{Kontaktplan}
	\acro{MA*}{Memory Bounded A*}
	\acro{OB}{Organisationsbaustein}
	\acro{PLC}{Programmable Logic Controller}
	\acro{RTA*}{Real-Time A*}
	\acro{SCL}{Structured Control Language}
	\acro{SPS}{Speicherprogrammierbare Steuerung}
	\acrodefplural{SPS}{Speicherprogrammierbare Steuerungen}
	\acro{STEP7}{STeuerungen Einfach Programmieren Version 7}
	\acro{TCP}{Transport Control Protocol}
	\acro{TIA-Portal}{Totally Integrated Automation Portal}
	\acro{UDP}{User Datagramm Protocol}

\end{acronym}

%%Nomenclature definitions


%	\nomenclature{CPS}{Cyber-physisches System}
%	\nomenclature{CPPS}{Cyber-physisches Produktionssystem}
%	%\acrodefplural{CPPS}{Cyber-physische Produktionssysteme}
%	\nomenclature{IoT}{Internet of Things}
%	\nomenclature{PLC}{Programmable Logic Controller}
%	\nomenclature{SPS}{Speicherprogrammierbare Steuerung}




\clearpage
\pagestyle{headings}
\pagenumbering{arabic}
  
	
	%% Hier kommt der eigentlich Text der Arbeit
	%% sinnvoll ist, für jedes Kapitel eine eigene 
	%% Datei zu nehmen, verkürzt die Zeit beim LaTeX-Durchlauf
	%% Erst am Schluss alle Dateien verwenden.
	
	%\chapter{Einführung in das Thema Industrie 4.0}

Der Begriff "`Industrie 4.0\index{Industrie 4.0}"' ist seit seiner Einführung durch die Bundeskanzlerin auf der Hannovermesse 2013 vor allem in den Bereichen Produktion und Fertigung in aller Munde. Da er jedoch je nach Branche unterschiedliche Bedeutungen haben kann, möchte ich als Einführung zunächst einmal erläutern, was Industrie 4.0 im Kontext der Automatisierungstechnik bedeutet.
%% Nachdem im Verlauf der vorhergehenden industriellen Revolutionen Menschliche Arbeitskraft zunächst durch Maschinen ersetzt wurde(erste industrielle Revolution), diese dann im Zuge der Massenproduktion zusammengeschaltet wurden(zweite industrielle Revolution) und dann Mitte diesen Jahrhunderts durch den Einsatz von elektrischen und informationstechnischen 
Der Begriff beschreibt die vierte industrielle Revolution und die damit einhergehende Verflechtung von informationstechnisch erhobenen Daten in den Produktionsablauf. Ein interessanter Aspekt ist hier beispielsweise der Bereich der prädiktiven Wartung, bei dem Anhand der Auswertung empirischer Daten mögliche Anlagenausfälle frühzeitig erkannt und behoben werden können.\\
Für den weiteren Verlauf dieser Arbeit ist vor allem auch ein sogenanntes \ac{CPPS} von Bedeutung. Diese bestehen aus der Verbindung von einzelnen, dezentralen Objekten wie beispielsweise Produktionsanlagen oder Logistikkomponenten, welche mit eingebetteten Systemen ausgestattet und zudem kommunikationsfähig gemacht werden. Durch eingebaute Sensoren und Aktoren kann die Umwelt erfasst und beeinflusst werden. Mittels der Kommunikationskomponenten können Daten aus der Produktion über ein Netzwerk oder das Internet ausgetauscht, beziehungsweise von entsprechenden Diensten ausgewertet, verarbeitet oder gespeichert werden.
\cite{Bauerhansl2014}\\
Sind mehrere \aclp{CPPS} an einem Produktionsprozess beteiligt, unabhängig von ihrem Standort, so spricht man auch von einer Smart Factory\index{Smart Factory}.
Abschließend ist noch zu erwähnen, dass der Begriff "`Industrie 4.0"' vor allem im deutschsprachigen Raum verwendet wird. Im internationalen Kontext werden viele der zentralen Punkte von Industrie 4.0 durch das Konzept des \ac{IoT} abgedeckt.




 
	
	
	%% Beim Anhang gibt es eine Musterdatei, die 
	%% ergänzt werden kann
	%\include{Anhang}
	
	
	%% Stichwortverzeichnis ausgeben.
	%% falls nicht notwendig, die folgende Zeile loeschen
	\printindex
	
\end{document}